\documentclass[10pt]{article}
\usepackage[margin=20mm]{geometry}
\usepackage{amsmath,amssymb,fontspec,unicode-math,xurl,hyperref}
\IfFontExistsTF{DejaVu Serif}{\setmainfont{DejaVu Serif}}{\setmainfont{Latin Modern Roman}}
\setmathfont{latinmodern-math.otf}
\setlength{\parindent}{0pt}\setlength{\parskip}{6pt}\sloppy
\begin{document}
Abstract: Four demonstrated defects affect the principal sparsification result: loss of hyperedge-labelled weight in the leverage construction, fixed-accuracy normalizer bias, an unjustified rank-dependent runtime simplification, and an invalid domain step in the chaining proof. Each is tied to a core claim or derivation and supports Major under the stated assumptions. The full two-term runtime is not contradicted by the simplification counterexample; the verified local repairs do not establish a complete corrected algorithm. Scientific adjudication remains pending.\par
Author: AI-assisted technical review\par
Date:\par
Title: Technical Verification of
Quantum Speedup for Hypergraph Sparsification\par
\# Identity, version and inspected scope\par
This report concerns Chenghua Liu, Minbo Gao, Zhengfeng Ji and Mingsheng Ying, Quantum Speedup for Hypergraph Sparsification, ICML 2025, PMLR 267, 38448–38466 [Retained latex: \textbackslash{}cite\{paper\}]; assignment 0111 / I25-43480. The reviewed file is the supplied 19-page published PDF. Page references below mean the printed PDF pages 1–19, not the proceedings page numbers. Its SHA-256 is recorded in the accompanying provenance and input-verification files.\par
The review covers the definitions, Algorithms 1–2, Theorems 3.4 and 4.1, their Appendices A–C, and the stated runtime dependencies in Section 5. The quantum primitives cited from other papers were treated as their stated interfaces, not re-proved or empirically implemented. Exact rational checks were independently written and run; no author code was run. The supplied manifest identifies no available author repository and records an OpenReview access gap for separate supplements. The embedded appendices were available and inspected.\par
\# Definitions and dependencies\par
For a hypergraph \(H=(V,E,w)\), \[Q_H(x)=\sum_{e\in E}w_e\max_{i,j\in e}(x_i-x_j)^2.\] A spectral sparsifier must approximate this energy multiplicatively for every \(x\). A sparse underlying graph assigns conductances \(c_{e,f}\) to a star \(S_e\) inside each hyperedge, with \(\sum_{f\in S_e}c_{e,f}=w_e\) (Definition 3.3, p. 6). A leverage-score overestimate must dominate \(w_e\max_{i,j\in e}R_{ij}\) for some underlying graph satisfying those groupwise constraints, and have total mass \(O(n)\) (Definition 3.2).\par
Algorithm 1 sparsifies each underlying graph, stores the resulting conductances, and uses them in the next weight update and the final overestimate. Algorithm 2 samples according to the resulting scores and estimates their sum before reweighting. Thus both interfaces are direct dependencies of Theorem 4.1; the cut applications inherit them.\par
\# Substantive findings\par
Problem 1 (Loss of hyperedge mass invalidates the leverage-score construction — Major).\par
\#\# Claim and exact location\par
Algorithm 1, lines 3–6 and 9 (p. 7), stores \(\widetilde G^{(t)}\) and passes that stored graph to WeightCompute and QOverestimate. Proposition B.5, equation (7) (p. 14), divides by \[\Delta_e=\sum_{g\in S_e}\widetilde c_{e,g}^{(t)}\widetilde R_g^{(t)}.\] Proposition B.7 (pp. 15–16) treats the averaged sparsified conductances as an underlying graph of \(H\) and telescopes ratios of successive sparsified conductances to establish domination.\par
\#\# Complete check\par
Take \(n=201\) vertices. Include every two-vertex hyperedge with weight one, and one additional hyperedge \(e_0=\{1,2,3\}\) of weight \(d=1/100\). This is a simple weighted hypergraph of rank three with \(m=20101\). Choose vertex 1 as the center of the star for \(e_0\). WeightInitialize produces the multigraph \[L_G=L_{K_n}+\frac d2 L_{\{\{1,2\},\{1,3\}\}}.\] All required weights are nonnegative, and the groupwise sums equal the original hyperedge weights.\par
Consider the graph-sparsifier output \(\widetilde G=K_n\), retaining the unit edges with their original labels and removing the two edges labelled by \(e_0\). On \(\mathbf1^\perp\), \(L_{K_n}=nI\) and the star Laplacian has largest eigenvalue three. Consequently, \[0\preceq L_G-L_{\widetilde G}\preceq
 \frac{3d/2}{n+3d/2}L_G=\frac1{13401}L_G.\] This output is therefore a valid \(0.1\)-spectral sparsifier. The input has 20102 underlying edges, so \(0.1\ge\sqrt{201/20102}\), as required by the stated graph-sparsification interface. The output has 20100 edges, exactly \(n/(0.1)^2\), and meets its size requirement.\par
Nevertheless, both stored conductances for \(e_0\) vanish. Equation (7) is \(0/0\) for that hyperedge. Moreover, already with \(T=1\), QOverestimate returns \(z_{e_0}=0\) regardless of its positive prefactor. Its purported underlying graph has group mass zero instead of \(1/100\). If that graph is interpreted as \(K_n\), its effective resistance between any two distinct vertices is \(2/n\), and \[w_{e_0}R_{e_0}=\frac1{10050}>0=z_{e_0}.\] This is not an issue about floating-point accuracy.\par
There are two further manifestations of the same failed invariant. A spectral approximation preserves aggregate quadratic forms, not the sum of conductances with each hyperedge label. Also, WeightCompute defines the pre-sparsification next conductances \(c^{(t+1)}\); equality (b) on p. 16 substitutes the independently reweighted \(\widetilde c^{(t+1)}\) into the recurrence. The proposed telescoping identity does not follow.\par
\#\# Counterevidence and bounded consequence\par
The example is an admissible output of the stated graph-sparsifier interface; it is not a measurement of the distribution produced by a specific implementation of that primitive. That distinction does not rescue the proof: no group-preservation property is supplied, and a sparsifier with fewer edges than positive-weight hyperedges cannot retain a labelled edge from every group. Retaining the unsparsified conductances could change the construction, but would require a new query-cost and domination argument. No such repaired algorithm is verified here.\par
The demonstrated failure directly affects Theorem 3.4, which supplies the core overestimates for Theorem 4.1. This supports Major even without claiming that quantum hypergraph sparsification is impossible or that every execution fails.\par
Problem 2 (Constant normalizer error survives sampling and invalidates the symmetrization — Major).\par
\#\# Claim and exact location\par
Algorithm 2, line 4 (p. 7), calls SumEstimate with accuracy \(0.1\), while line 8 uses the returned sum in the output weights. Appendix C, pp. 17–19, explicitly retains this fixed accuracy and claims an arbitrarily small \(\varepsilon\)-error after choosing a sufficiently large sample size. The equality introducing Rademacher signs in the uncentered second term on p. 18 is invalid.\par
\#\# Derivation and an admissible family\par
Grant an exact valid score oracle, so this finding is independent of Problem 1. Put \(s=\sum_e z_e\), \(p_e=z_e/s\), and let \(\widehat s=(1+\delta)s\). Conditional on this normalizer, the returned energy is \[\widehat Q(x)=\frac1M\sum_{t=1}^M
       \frac{\widehat s\,w_{e_t}}{z_{e_t}}Q_{e_t}(x),
 \qquad \mathbb E[\widehat Q(x)\mid\widehat s]=(1+\delta)Q_H(x).\] Jensen’s inequality already gives \(\mathbb E|\widehat Q(x)-Q_H(x)|\ge|\delta|Q_H(x)\). Increasing \(M\) cannot remove this bias.\par
A family makes the failure deterministic conditional on an allowed normalizer value. Let \[V=\{0,\ldots,n-1\},\quad
 E=\{\{0,i,j\}:1\le i<j\le n-1\},\quad w_e=1.\] Every hyperedge has rank three. Use the star centered at 0 with conductance \(1/2\) on each of its two edges. The resulting graph is a star whose leaf conductances are \((n-2)/2\). Hence each hyperedge has maximum resistance \(4/(n-2)\). Set \(z_e=4/(n-2)\), a valid overestimate with \[m=\binom{n-1}{2},\qquad s=2(n-1),\qquad p_e=1/m.\] For \(x_0=1\) and all other coordinates zero, every \(Q_e(x)=1\) and \(Q_H(x)=m\). Every possible sampled sequence yields \[\widehat Q(x)=\frac{\widehat s}{s}m.\] The normalizer \(\widehat s=21s/20\) satisfies the stated 10\% accuracy guarantee but produces relative error \(1/20\), exceeding \(\varepsilon=1/100\) for every sample sequence. For any fixed hidden constant in \(M=\Theta(n\log n\log3/\varepsilon^2)\), choose \(n\) sufficiently large that \(M\le\binom{n-1}{2}\); such an integer exists since \(n/\log n\to\infty\). Then the sampling primitive’s \(M\le m\) condition and the sparsification regime \(\varepsilon\ge\sqrt{n/m}\) are both met. No single-edge example outside that regime is needed.\par
On p. 18, the uncentered mean term is replaced by a Rademacher-signed sum. Those quantities do not have the same expectation: in the displayed family the former is a nonzero constant while the latter decreases with sample size. The preceding triangle-inequality coefficient also needs \(|s/\widehat s-1|\), rather than a signed coefficient.\par
\#\# Counterevidence and correction bound\par
A particular normalizer implementation could return a much more accurate value on a particular input. The result above disproves the claimed consequence of the stated \(0.1\) interface and the displayed proof, not a universal lower bound on all implementations.\par
A defensible partial correction is to estimate \(s\) to \(\eta=O(\varepsilon)\) and allocate error to the unbiased sampling step. If an unbiased-sum estimator has relative error at most \(\beta\), the rescaled estimator has error at most \[\eta+(1+\eta)\beta.\] For \(0<\varepsilon\le1\), taking \(\eta=\beta=\varepsilon/3\) makes this at most \(7\varepsilon/9<\varepsilon\). The additional stated sum-estimation cost is \(\widetilde O(r\sqrt m/\varepsilon)\), within the sampling term. This verifies the normalization repair only; it does not fix Problem 1 or prove the required simultaneous high-probability sampling bound. Appendix C ends with an expectation bound, so that probability upgrade remains unverified.\par
Problem 3 (The density conditions do not justify the advertised runtime simplification — Major).\par
Remark 4.2 (p. 7) drops \(r\sqrt{mnr}\) under \(\varepsilon\ge\sqrt{n/m}\) and \(m\ge nr\). The exact ratio of this term to \(r\sqrt{mn}/\varepsilon\) is \(\varepsilon\sqrt r\); density cancels. With rank growing polynomially, fixed \(\varepsilon\) and a sufficiently dense simple hypergraph, that ratio need not be polylogarithmic. For example, include all three-vertex hyperedges and the full vertex set, with \(n>3\). Then \(r=n\) and \(m=\binom n3+1=\Theta(n^3)\); both stated inequalities hold for sufficiently large \(n\) and \(\varepsilon=1/2\). The ratio is \(\sqrt n/2\), which is not polylogarithmic in any of the stated parameters.\par
The abstract (p. 1) advertises the simplified bound as the main algorithmic runtime for rank \(r\), and Remark 4.2 supplies its stated justification. Thus the lost nonpolylogarithmic factor affects a core complexity claim and supports Major. The full two-term complexity in Theorem 4.1 and Table 1 is not contradicted by this arithmetic. Retain both terms in the advertised bound or explicitly restrict, for example, to \(\varepsilon\le1/\sqrt r\) (up to suppressed logarithms). The validity of the unsimplified bound and the availability of this local correction do not validate the stronger claim.\par
Problem 4 (Proposition C.2 leaves a core chaining hypothesis unjustified — Major).\par
Proposition C.2 (p. 17) asserts \(Q_H(L_G^{+/2}x)\ge\|x\|^2\) for every \(x\perp\mathbf1\). For the graph with one unit edge \(\{1,2\}\) and isolated vertex 3, take \(x=(1,1,-2)\). This satisfies \(x\perp\mathbf1\) but lies in \(\ker L_G\), giving left side zero and right side six.\par
The proof establishes only \[Q_H(L_G^{+/2}x)\ge x^\top P_{\operatorname{range}L_G}x.\] Thus \(x\in(\ker L_G)^\perp\) is a sufficient corrected assumption. Even for a connected graph, the set \(T\) on p. 18 as literally defined contains arbitrary kernel directions and is not contained in the Euclidean unit ball. Restricting the chaining domain to the range leaves the represented energy unchanged whenever the underlying graph has the same nullspace as \(H\). The proof on p. 19 explicitly invokes Proposition C.2 to assert \(T\subseteq B_2^n\), the hypothesis needed to apply Lemma C.1 in the proof of Theorem 4.1. This false statement and its invalid application therefore affect a core derivation and support Major. The range restriction verifies a local domain correction, not the remaining chaining argument. Disconnected inputs are allowed by the paper; connectivity must not be silently imposed.\par
\# Overall impact, open scope and conclusion\par
Problems 1, 2 and 4 invalidate steps directly supporting the main algorithmic guarantee. Problem 3 independently invalidates the advertised rank-dependent runtime simplification. Each supports Major through its stated core dependency; the shortness of a local repair does not change the impact of the original claim or derivation. The general sparsification objective, imported quantum primitives, and possibility of a corrected algorithm remain open within this review.\par
The Section 5 substitutions also retain unresolved accuracy dependence: inserting \(M=\widetilde O(n/\varepsilon^2)\) into the stated classical mincut and flow runtimes yields respectively \(\widetilde O(rn^2/\varepsilon^2)\) and \(\widetilde O(rn^{3/2}/\varepsilon^3)\) under direct composition. No alternative accuracy-independent downstream algorithm was verified. This is recorded as Needs Evidence for those corollaries, rather than a claim that the asserted algorithms cannot exist.\par
A bounded search of the official proceedings page and arXiv submission history found the latter listing v1 (3 May 2025); no later correction was identified in the material consulted [Retained latex: \textbackslash{}cite\{record\}]. This is not an exhaustive errata search, and no author adoption is claimed.\par
Dimension | State\par
Supported issue judgment | Major: Problems 1–4\par
Scientific adjudication | Pending; initial review only\par
Repair verification | Normalizer error budget and range restriction only\par
Author adoption | Not requested; unknown\par
Remaining work | Correct leverage oracle, full sampling probability bound, downstream accuracy costs\par
Human review | None recorded\par
9 C. Liu, M. Gao, Z. Ji and M. Ying, Quantum Speedup for Hypergraph Sparsification, PMLR 267 (2025), 38448–38466. Published PDF supplied for this assignment. https://proceedings.mlr.press/v267/liu25r.html. The same authors, arXiv:2505.01763, submission-history record listing v1, 3 May 2025, accessed 24 October 2026. Metadata only; a second full-text version was not substituted. https://arxiv.org/abs/2505.01763.\par
\# Evidence and reproducibility\par
Evidence under `verification/0111/evidence/` includes the extracted original text, the rendered source page 18, `independent\_checks.py` and its actual JSON output, and source/hash verification. The rational script confirms \(1/13401\), \(1/10050\), the normalizer discrepancy and the disconnected-vector calculation. The general family and error-budget arguments above are analytic proofs; the script is an arithmetic check, not an empirical quantum-algorithm reproduction.\par
The self-contained TeX was compiled with an existing pdfLaTeX installation in the provided execution runtime; its exact bytes, PDF hash, build logs and visual-inspection record are bound in provenance. The final files are saved in the assigned Windows staging tree. All scientific analysis was performed in this conversation. A support agent only located tools and skills; it did not review papers. Scientific decisions remain pending, and the human-review list is empty.\par
\end{document}
